\documentclass[10pt,a4paper]{article}
\usepackage[utf8]{inputenc}
\usepackage[T1]{fontenc}
\usepackage[french]{babel}
\usepackage{amsmath, amssymb}
\usepackage{geometry}
\usepackage{graphicx}
\geometry{hmargin=2cm,vmargin=2cm}

\title{Entraînement : Automatismes 01}
\author{Partie 1 - 6 points - 20 minutes}
\maDate{}

\begin{document}

\maketitle

\textit{Consigne : Pour chaque question, recopier sur la copie son numéro et la réponse correspondante. Aucune justification n'est demandée.} 

\begin{center}
\renewcommand{\arraystretch}{2.5}
\begin{tabular}{|p{0.1\textwidth}|p{0.5\textwidth}|p{0.35\textwidth}|}
\hline
\textbf{N°} & \textbf{Énoncé} & \textbf{Réponse} \\ \hline
1 & Quel est le quart de 32 ? &  \\ \hline
2 & Un trajet dure 150 min. Quelle est sa durée en heures et minutes ? &  \\ \hline
3 & Voici une série de notes : $10 \, ; 7 \, ; 15 \, ; 12 \, ; 9$. Quelle est la médiane ? &  \\ \hline
4 & Dans un triangle $DEF$ rectangle en $E$, on sait que $\widehat{D} = 42^\circ$. Calculer la mesure de l'angle $\widehat{F}$. &  \\ \hline
5 & Sur une droite graduée, l'abscisse de $A$ est $\frac{3}{4}$ et celle de $B$ est $\frac{3}{2}$. Quel point est le plus éloigné de l'origine ? &  \\ \hline
6 & Dans le triangle $ABC$ rectangle en $C$, quel est le cosinus de l'angle $\widehat{ABC}$ ? (Donner la formule avec les côtés) &  \\ \hline
7 & Un article coûte 60 €. On applique une réduction de 10 \%. Quel est le nouveau prix ? &  \\ \hline
8 & Dans un collège de 400 élèves, 30 \% sont externes. Combien d'élèves sont externes ? &  \\ \hline
9 & Compléter le script Scratch pour dessiner un triangle équilatéral : \newline \textbf{répéter \dots\dots fois} \newline \textbf{avancer de 50} \newline \textbf{tourner de \dots\dots degrés} &  \\ \hline
10 & Si $AD = 6$ cm et que le rapport de réduction est de 0,5, quelle est la longueur de l'image $A'D'$ ? &  \\ \hline
\end{tabular}
\end{center}

\newpage

\section*{Correction (pour s'auto-évaluer)}

\begin{enumerate}
    \item \textbf{8} (car $32 \div 4 = 8$). 
    \item \textbf{2 h 30 min} (car $150 = 2 \times 60 + 30$). 
    \item \textbf{10} (Série ordonnée : $7 ; 9 ; \mathbf{10} ; 12 ; 15$). 
    \item \textbf{$48^\circ$} (car $90 - 42 = 48$). 
    \item \textbf{Le point $B$} (car $\frac{3}{2} = 1,5$ et $\frac{3}{4} = 0,75$). 
    \item \textbf{$\cos(\widehat{ABC}) = \frac{BC}{AB}$} (Côté adjacent / Hypoténuse).
    \item \textbf{54 €} (La réduction est de 6 €). 
    \item \textbf{120 élèves} (car $0,30 \times 400 = 120$). 
    \item \textbf{Répéter 3 fois} et \textbf{Tourner de 120 degrés}. 
    \item \textbf{3 cm} (car $6 \times 0,5 = 3$). 
\end{enumerate}

\end{document}